\documentclass[10pt,a4paper]{article}
\usepackage[margin=1.8cm]{geometry}
\usepackage{graphicx,booktabs,tabularx,amsmath,xcolor,hyperref,caption,subcaption,enumitem}
\captionsetup{font=small,labelfont=bf,skip=3pt}
\setlist{nosep,leftmargin=*}
\setlength{\parskip}{2pt}
\hypersetup{colorlinks,linkcolor=black,citecolor=black,urlcolor=blue!50!black}
\newcommand{\fig}[1]{../results/figures/#1}
\newcommand{\viz}[1]{../results/#1}
\newcommand{\fillin}[1]{\colorbox{yellow!40}{\textbf{[#1]}}}
\title{\vspace{-1.2cm}\textbf{Network Dissection of AlexNet and VGG16:\\ a PyTorch re-implementation and an evaluation of the method}}
\author{Interpretable Machine Learning --- Group Project \\ \fillin{Group members (names, roll numbers)}}
\date{\fillin{Submission date} \quad Code: \fillin{GitHub link / zip name}}
\begin{document}
\maketitle
\vspace{-0.6cm}
\begin{abstract}\noindent
Network Dissection \cite{bau2017} measures how interpretable the hidden units of a CNN are by scoring every
unit as a segmenter of human-labelled concepts. We re-implemented the full method in PyTorch, dissected
AlexNet (ImageNet, Places365, random) and VGG16 (with and without batch normalisation) on the Broden dataset,
and repeated every experiment of the paper (Sec.~3.1--3.7), the training experiments at reduced scale. Most
findings are reproduced: concepts move from colours/textures in early layers to objects in conv5, VGG16 is
more interpretable than AlexNet, a random rotation of conv5 destroys about 69\% of the unique detectors without
changing accuracy, and interpretability correlates with transfer accuracy. Some claims are \emph{not}
reproduced: batch normalisation does not reduce interpretability in our models, removing dropout has no clear
effect, and the Places365 $>$ ImageNet ordering is small and flips with the IoU threshold. Using the evaluation criteria of Molnar's book
\cite{molnar} we further show that unit labels have low single-unit fidelity but meaningful group fidelity,
are moderately stable, pass a randomisation sanity check, and that learned concept directions beat single units.
\end{abstract}

\section{Introduction}
In Molnar's taxonomy \cite{molnar} Network Dissection is a \emph{post-hoc, model-specific, global} method from
the chapter on learned features: it answers ``what concept does this unit detect?'' for every unit of a trained
CNN. The paper tests whether interpretable units are an accident of the coordinate system (H1: any direction is
as interpretable as a unit) or whether training finds a special, \emph{axis-aligned} basis (H2), and how
interpretability depends on layer, architecture, training data and training tricks. Our contributions:
(i) a from-scratch PyTorch implementation with exact quantiles and receptive-field-anchored upsampling, tested
against a brute-force IoU; (ii) a reproduction of each experiment of the paper for AlexNet and VGG16;
(iii) an evaluation of the method itself with the book's criteria (fidelity, stability, sanity checks,
comparison with other methods).

\section{Datasets and why we used them}
\textbf{Broden} (\texttt{broden1\_224}) is the dataset the method is defined on. It merges ADE20K, Pascal-Context,
Pascal-Part, OpenSurfaces and DTD plus generated colour labels: 63{,}305 images (224\,px) with labels at
112\,px for 1{,}197 concepts in six categories --- 584 objects, 234 parts, 468 scenes, 32 materials, 47 textures,
11 colours (exactly the paper's Table~1). Objects, parts, materials and colours are labelled per pixel,
scenes and textures per image. Pixel-level labels are what dissection needs (a unit's activation map is
scored as a segmentation), and using the paper's own data makes our numbers comparable. The official server
was offline, so we used the Internet Archive copy and verified it byte-for-byte (SHA-1 of the archived
original). \textbf{Imagenette} (10 ImageNet classes, 9{,}469 train / 3{,}925 val images) was used where class labels
of the network's own task are needed: the accuracy check of the rotation experiment, the fidelity test, and
training small AlexNets for the paper's training-condition experiments. \fillin{Mention instructor approval of Imagenette if obtained}

\section{Libraries}
\begin{tabularx}{\linewidth}{@{}lX@{}}\toprule
PyTorch 2.7 & tensors, GPU computation, autograd (training, feature visualisation, Net2Vec) \\
torchvision 0.22 & AlexNet/VGG16/VGG16-BN definitions and ImageNet weights; Imagenette loading and transforms \\
NumPy, pandas & memory-mapped Broden arrays; reading Broden's CSV files; result tables \\
Pillow & decoding JPEGs and the RG-encoded label PNGs; writing image panels \\
SciPy & real Schur decomposition (fractional rotations); statistical tests (Mann--Whitney, Wilcoxon, Spearman) \\
Matplotlib & figures \\\bottomrule
\end{tabularx}
No implementation of Network Dissection, Net2Vec, feature visualisation or TCAV was imported; these are written
in the repository ($\approx$1{,}850 lines including tests).

\section{Method and implementation}
For a layer with $K$ units and Broden images $x$:
(1)~$T_k$ is the top quantile with $P(a_k>T_k)=0.005$ over all spatial positions of all images;
(2)~the activation map $A_k(x)$ is bilinearly upsampled to $112\times112$ with each activation placed at the centre of
its receptive field, giving $S_k(x)$, and binarised $M_k(x)=S_k(x)>T_k$;
(3)~for every concept $c$ with label mask $L_c(x)$,
$\mathrm{IoU}_{k,c}=\sum_x|M_k(x)\cap L_c(x)|\,/\,\sum_x|M_k(x)\cup L_c(x)|$, where the sums run only over images
that have a label of $c$'s category; (4)~unit $k$ is a \emph{detector} of its best concept if the IoU $>0.04$, and
the interpretability of a layer is its number of \textbf{unique detectors} (distinct concepts).

\textbf{Implementation.} Intersections for all 1{,}197 concepts are accumulated per batch with a GPU
\texttt{index\_add} into a concept$\times$unit table (exact in float32 because per-batch counts are $<2^{24}$);
scene/texture labels act as full-image masks, so their intersection is the mask area; the union is
$|M|+|L|-|M\cap L|$ with $|M|$ restricted per category. The quantile is \emph{exact}: a subsample proposes cut
levels, one streaming pass counts values above/equal to each level, a second pass keeps only values above the
chosen level and takes the order statistic; ties (units that are 0 on $>$99.5\% of positions) give $T_k=0$.
All sums are stored per image shard (50 random shards), so bootstrap intervals need no recomputation.
Receptive-field centres are computed from the conv/pool geometry (AlexNet conv1: start 3, stride 4).
A random rotation $Q\in SO(256)$ is drawn by QR of a Gaussian matrix; $Q^\alpha$ is computed from the real Schur
form by multiplying every rotation angle by $\alpha$. One AlexNet (5 layers) is dissected in $\approx$9 minutes.

\section{Reproduction of the paper's experiments}
\textbf{Differences from the paper} (to read the numbers correctly): torchvision AlexNet (64 conv1 filters, no
LRN, 224\,px) instead of Caffe AlexNet (96 filters, LRN, 227\,px); the CSAIL PyTorch \emph{Places365} AlexNet instead
of Places205; training experiments on Imagenette instead of Places205; transfer task Broden scenes instead of
Action40. GoogLeNet, ResNet and self-supervised models are outside our topic (AlexNet/VGG).

\subsection{Human evaluation (paper Sec.~3.1, Table~3)}
Protocol: 30 randomly chosen detector units of AlexNet-Places365 conv5; each shown as its 15 most activating
Broden images with the unit mask highlighted (as in the paper), together with a phrase. Unlike the paper we show
the true label for a random half of the units and a decoy label of the same category for the other half, so that
raters' tendency to say ``yes'' can be measured. Raters work independently.
\begin{center}\small\begin{tabular}{lccc}\toprule
& true label judged descriptive & decoy judged descriptive & Fleiss' $\kappa$ \\\midrule
Paper, conv5 (AlexNet-Places205) & 71\% & -- & -- \\
Ours (\fillin{n} raters, 30 units) & \fillin{\%} & \fillin{\%} & \fillin{value} \\\bottomrule
\end{tabular}\end{center}
\noindent\fillin{Run the study (visualize.py, ~20 min) and fill this table; otherwise remove this subsection.}

\subsection{Is interpretability axis-aligned? (Sec.~3.2, Fig.~4)}
\begin{figure}[t]\centering
\begin{subfigure}{0.49\linewidth}\includegraphics[width=\linewidth]{\fig{fig_rot-alexnet-places365.png}}\caption{Rotation $Q^\alpha$ of AlexNet-Places365 conv5 (3 seeds).}\label{fig:rot}\end{subfigure}\hfill
\begin{subfigure}{0.49\linewidth}\includegraphics[width=\linewidth]{\fig{fig_architectures.png}}\caption{Last conv layer, bootstrap 95\% intervals.}\label{fig:arch}\end{subfigure}
\caption{Left: interpretability drops as the basis is rotated away from the learned one (paper Fig.~4). Right: architecture, training data and batch norm (paper Fig.~7).}
\end{figure}
Fig.~\ref{fig:rot}: rotating conv5 of AlexNet-Places365 reduces unique detectors from 55 to $17.3\pm1.7$ at
$\alpha=1$ (three random rotations), a \textbf{69\% drop}; for AlexNet-ImageNet 53$\to$23 (57\%, one seed). The paper
reports 80\%. As in the paper the decrease is gradual and strongest between $\alpha=0.6$ and $0.8$. The paper's premise
that a rotation keeps discriminative power was \emph{measured}: feeding $Q^\top Q f(x)$ to the rest of the network gives
identical Imagenette accuracy (98.4\%, max logit difference $1.5\cdot10^{-5}$), while $Qf(x)$ without the inverse
gives 9.6\% (chance). Interpretability is therefore a property of the basis, not of the information: \textbf{H2 is supported}.

\subsection{Concepts by layer (Sec.~3.3, Fig.~5)}
\begin{figure}[h]\centering\includegraphics[width=0.86\linewidth]{\fig{fig_layers_alexnet.png}}
\caption{Unique detectors per layer of AlexNet (paper Fig.~5).}\label{fig:layers}\end{figure}
Fig.~\ref{fig:layers}: colour and texture detectors dominate conv1--conv2 and object/part detectors appear mainly in
conv5 (AlexNet-ImageNet: 6, 20, 33, 28, 53 unique detectors; objects 0$\to$17). The count is not monotone
(conv3 $>$ conv4), but conv3 has 384 units and conv4 256; per 100 units they are equal (8.6 vs 10.9).
Places365 has \emph{fewer} detectors than ImageNet in conv2--conv3 (21 vs 33 at conv3): in this checkpoint 237 of
384 conv3 units are zero on more than 99.5\% of positions (we verified the checkpoint's scene predictions on Broden
are sensible, e.g.\ corridor$\to$corridor, so this is a property of the model, not a loading error).

\subsection{Architectures and training data (Sec.~3.4, Fig.~7)}
\begin{figure}[t]\centering
\begin{subfigure}{0.49\linewidth}\includegraphics[width=\linewidth]{\fig{fig_training_variants.png}}\caption{AlexNets trained on Imagenette (variant, val.\ accuracy).}\label{fig:train}\end{subfigure}\hfill
\begin{subfigure}{0.46\linewidth}\includegraphics[width=\linewidth]{\fig{fig_training_dynamics.png}}\caption{Baseline during training.}\label{fig:dyn}\end{subfigure}
\caption{Training conditions at reduced scale (paper Fig.~10, 11, 13).}
\end{figure}
Fig.~\ref{fig:arch}: \textbf{VGG16 $>$ AlexNet} (89 vs 53 unique detectors, 37 vs 17 objects), as in the paper. A random
AlexNet has only 6 (paper: few), mostly the texture ``chequered'' that random filters pick up. \textbf{Places365 vs ImageNet}:
55 vs 53 with overlapping intervals ([52, 60] vs [49, 57]); Places365 has more object (20 vs 17) and scene (5 vs 1)
detectors, the direction the paper reports, but the difference is not significant and it reverses for IoU thresholds
$\geq0.06$ (Sec.~\ref{sec:stab}), contradicting the paper's statement that relative orderings are threshold-stable.

\subsection{Training conditions (Sec.~3.5, Fig.~10--11) and width (Sec.~3.7, Fig.~13)}
We trained six AlexNets from scratch on Imagenette with one recipe (60 epochs, SGD, one-cycle LR) and dissected conv5
(Fig.~\ref{fig:train}). \emph{Random initialisation:} three seeds give 17, 20, 17 unique detectors at the same accuracy
($\approx$83\%) --- reproduced. \emph{Training iterations:} 6 detectors at initialisation, 6--8 up to iteration 400,
15 at 800--1600 and 17 at the end; object detectors appear only after $\approx$800 iterations --- reproduced in shape (the paper sees objects after $\approx$10{,}000 iterations on a 250$\times$ larger dataset).
\emph{Dropout:} the paper finds more textures and fewer objects without dropout; we see 16 vs 17 unique and
objects 5 vs 6 --- \textbf{not reproduced} (no clear effect). \emph{Batch norm:} the paper finds a strong decrease; our BN model
has \textbf{more} detectors (26, interval [23, 29], vs 17, [15, 21]) and higher accuracy (85.9 vs 82.6\%) --- \textbf{not reproduced}.
For pretrained VGG16, BN gives fewer detectors in conv3\_3 (2 vs 9) and conv4\_3 (26 vs 46) but slightly more in
conv5\_3 (96 vs 89); these two checkpoints were trained separately, so this comparison is confounded.
\emph{Width:} widening conv5 to 768 units with global average pooling gives many more detector units (300 vs 56)
but not more unique concepts (14 vs 17) --- the paper reports both increasing, so only half is reproduced.

\subsection{Discrimination vs interpretability (Sec.~3.6, Fig.~12)}
\begin{figure}[t]\centering
\begin{subfigure}{0.49\linewidth}\includegraphics[width=\linewidth]{\fig{fig_discrimination.png}}\caption{Transfer accuracy vs unique object detectors.}\label{fig:disc}\end{subfigure}\hfill
\begin{subfigure}{0.49\linewidth}\includegraphics[width=\linewidth]{\fig{fig_cascade.png}}\caption{Cascading randomisation sanity check.}\label{fig:casc}\end{subfigure}
\caption{Left: paper Fig.~12 (our transfer task: 30 Broden scenes). Right: re-initialising layers from the top.}
\end{figure}
A linear classifier on pooled last-conv features (30 most frequent Broden scenes, 6{,}844 train / 2{,}789 val images)
against the number of unique object detectors (Fig.~\ref{fig:disc}) gives Spearman $\rho=0.79$ ($p=0.004$, 11 models),
the positive relation of the paper. Most of it comes from the gap between pretrained and small models: within the five
pretrained/random models the order is perfect ($\rho=1$), within our six small models there is none ($\rho=-0.38$, $p=0.45$).

\begin{table}[h]\centering\small
\begin{tabularx}{\linewidth}{@{}Xll@{}}\toprule
Claim of the paper & Our result & Verdict \\\midrule
Interpretability is axis-aligned; rotation keeps accuracy (3.2) & $-69\%$ unique detectors, accuracy identical & reproduced (smaller drop) \\
Low layers colour/texture, high layers objects/parts (3.3) & yes, objects 0$\to$17 from conv1 to conv5 & reproduced \\
VGG $>$ AlexNet; trained $\gg$ random (3.4) & 89 vs 53; 6 for random & reproduced \\
Places $>$ ImageNet (3.4) & 55 vs 53, intervals overlap, flips with threshold & weak / not significant \\
Similar interpretability across seeds (3.5) & 17, 20, 17 & reproduced \\
Detectors emerge with training (3.5) & 6 $\to$ 17, objects after $\sim$800 it. & reproduced \\
No dropout: more textures, fewer objects (3.5) & no clear difference & not reproduced \\
Batch norm decreases interpretability (3.5) & more detectors (small models); mixed (VGG) & not reproduced \\
More discriminative $\Rightarrow$ more detectors (3.6) & $\rho=0.79$ overall, driven by model groups & reproduced (weakly) \\
Wider layer $\Rightarrow$ more detectors (3.7) & more detector units, not more unique concepts & partly \\
Orderings stable across IoU thresholds (2.2) & Places/ImageNet order flips at $\geq0.06$ & not reproduced \\\bottomrule
\end{tabularx}
\caption{Summary of the reproduction.}\label{tab:summary}\end{table}

\section{Evaluating Network Dissection with the book's criteria}
\textbf{Fidelity (permutation importance).} Using Molnar's permutation feature importance with units as features,
we shuffled one conv5 unit at a time across Imagenette images (5 repeats). Units with a concept label are \emph{not} more
important than unlabelled ones (mean loss ratio 1.014 vs 1.015, Mann--Whitney $p=0.80$; VGG16 $p=0.63$) and IoU does not
correlate with importance (Spearman $-0.01$); the largest single-unit accuracy drop is 0.2\%. Because many units carry the
same concept, single-unit importance underestimates it (the book's warning about correlated features), so we also shuffled
all units with the same label together. Our one test fixed in advance --- the ``dog'' units for the dog class (English springer)
--- gives a 1000-way accuracy drop of \textbf{24.3 points} for AlexNet (13 units; 10{,}000 random 13-unit groups: mean 3.4, max 22.3,
$p=0.0001$, $z=7.2$) and \textbf{29.3 points} for VGG16 (35 units, $p=0.0017$, $z=5.7$). The 10-way drop is below 1 point: the dog units
carry \emph{which breed}, not \emph{dog vs.\ not dog}. So labels are faithful for groups of units, not for single units.

\textbf{Sanity check.} Adebayo et al.'s cascading randomisation (book, saliency chapter) applied to dissection
(Fig.~\ref{fig:casc}): re-initialising conv5 and above drops conv5 from 53 to 14 unique detectors, re-initialising
everything to 3. Re-initialising only the classifier changes \emph{nothing} (53, identical labels): dissection
describes what a layer represents, not how the network uses it --- which is why the fidelity test is needed.

\textbf{Stability.}\label{sec:stab} Bootstrapping the 50 image shards, 89--94\% of conv5 detector labels survive resampling
(bootstrap intervals are approximate: unique counts are biased upward under resampling, and for Places365 conv4 the observed
value lies just below its interval). Changing the activation quantile from 0.5\% to 1\% keeps 80--88\% of labels, 0.1\% only
57--61\% (unique counts 22--27 instead of 53--55); Gaussian pixel noise ($\sigma=0.05$/$0.1$) keeps 84--96\% / 70--73\%.
\begin{figure}[t]\centering
\begin{subfigure}{0.40\linewidth}\includegraphics[width=\linewidth]{\fig{fig_iou_threshold.png}}\caption{IoU threshold sweep.}\label{fig:thr}\end{subfigure}\hfill
\begin{subfigure}{0.52\linewidth}\includegraphics[width=\linewidth]{\fig{fig_negative.png}}\caption{Top vs bottom 0.5\% of activations.}\label{fig:neg}\end{subfigure}
\caption{Left: the number of unique detectors and the Places/ImageNet order depend on the 0.04 threshold. Right: negative activations.}
\end{figure}
The detector threshold matters a lot (Fig.~\ref{fig:thr}): 81 unique detectors at 0.02, 53 at 0.04, 22 at 0.08, and the
Places/ImageNet order flips. The 0.04 threshold is also weak: 22\% of the units of a \emph{random} AlexNet pass it.

\textbf{Negative activations.} The book notes that dissection ignores negative activations and suggests the lower quantile.
Dissecting the bottom 0.5\% of pre-ReLU activations (Fig.~\ref{fig:neg}) gives only 12 (AlexNet-ImageNet), 11 (Places365) and 16
(VGG16) unique concepts, almost all repetitive high-contrast textures (dotted, grid, chequered), and no unit has the same concept
for both tails. The negative tail is mostly not concept-specific.

\textbf{Units vs.\ directions} (book: concept chapter). We learned a linear direction in conv5 for 48--49 Broden concepts (Broden
train split, pixel-wise BCE; Net2Vec \cite{net2vec}) and dissected each direction exactly like a unit on the val split. Directions
beat the best single unit for 47/48 (ImageNet) and 49/49 (Places365) concepts (mean IoU 0.093 vs 0.063; Wilcoxon $p<10^{-12}$),
and capture 12--13 concepts that have no detector unit (building, floor, torso, \dots). With the rotation result this gives the
full picture: \emph{random} directions are less interpretable than units, \emph{learned} ones more --- concepts are partly
distributed over units.

\textbf{Prototypes and criticisms} (book: example-based explanations). Fig.~\ref{fig:units} shows each unit's most activating images
(prototypes), the top-100 images where its label fits least (criticisms) and an activation-maximisation image. The ``wheel'' unit
(IoU 0.18) also fires on plates and spirals, the ``dog'' unit on cats and sheep --- a concrete case of the ``illusion of
interpretability''. For texture labels only 3--47\% (median 18\%) of a unit's top-100 images carry the label (object/part labels: 30--97\%), because texture labels exist only
on the 5{,}640 DTD images and the IoU is computed only over those images; many ``texture detectors'' are partly a dataset artefact.
\begin{figure}[t]\centering\setlength{\tabcolsep}{1pt}\small
\begin{tabular}{lccc}
& prototypes (8 most activating images) & criticisms & feature vis. \\
\rotatebox{90}{\scriptsize wheel 0.18} & \includegraphics[height=1.15cm]{\viz{alexnet-imagenet/viz_conv5/unit0248_prototypes.jpg}} & \includegraphics[height=1.15cm]{\viz{alexnet-imagenet/viz_conv5/unit0248_criticisms.jpg}} & \includegraphics[height=1.15cm]{\viz{alexnet-imagenet/featviz/unit0248.png}} \\
\rotatebox{90}{\scriptsize dog 0.13} & \includegraphics[height=1.15cm]{\viz{alexnet-imagenet/viz_conv5/unit0176_prototypes.jpg}} & \includegraphics[height=1.15cm]{\viz{alexnet-imagenet/viz_conv5/unit0176_criticisms.jpg}} & \includegraphics[height=1.15cm]{\viz{alexnet-imagenet/featviz/unit0176.png}} \\
\rotatebox{90}{\scriptsize cat 0.17} & \includegraphics[height=1.15cm]{\viz{alexnet-imagenet/viz_conv5/unit0130_prototypes.jpg}} & \includegraphics[height=1.15cm]{\viz{alexnet-imagenet/viz_conv5/unit0130_criticisms.jpg}} & \includegraphics[height=1.15cm]{\viz{alexnet-imagenet/featviz/unit0130.png}} \\
\rotatebox{90}{\scriptsize grass 0.13} & \includegraphics[height=1.15cm]{\viz{alexnet-imagenet/viz_conv5/unit0095_prototypes.jpg}} & \includegraphics[height=1.15cm]{\viz{alexnet-imagenet/viz_conv5/unit0095_criticisms.jpg}} & \includegraphics[height=1.15cm]{\viz{alexnet-imagenet/featviz/unit0095.png}} \\
\end{tabular}
\caption{AlexNet-ImageNet conv5 units (label, IoU). Yellow borders: unit mask $M_k$. Criticisms of the dog unit are cats and sheep.}\label{fig:units}
\end{figure}

\textbf{Consistency.} The sets of concepts found in conv5 overlap with Jaccard 0.61 (AlexNet ImageNet vs Places365) and 0.46
(AlexNet vs VGG16, both ImageNet); random AlexNet shares almost nothing (0.02--0.09). Different models trained on similar data
use similar, but far from identical, concepts (the Rashomon effect of the book).

\begin{table}[h]\centering\small
\begin{tabularx}{\linewidth}{@{}lX@{}}\toprule
Property \cite{molnar} & Network Dissection \\\midrule
Expressive power & one concept word per unit, from a fixed vocabulary (Broden) \\
Translucency / portability & high translucency (reads activations); only CNNs and only concepts that Broden labels \\
Algorithmic complexity & one forward pass + IoU table; $\approx$2 min per AlexNet layer on one GPU \\
Fidelity & low for single units, high for groups of units with the same label \\
Stability & 89--94\% of conv5 labels under resampling; sensitive to the quantile and IoU thresholds \\
Comprehensibility & \fillin{from the human study} \\
Degree of importance & not provided by IoU --- needs a separate test (permutation importance) \\\bottomrule
\end{tabularx}
\caption{Network Dissection judged by the properties of explanation methods in the book.}\end{table}

\section{Challenges and how we solved them}
\begin{tabularx}{\linewidth}{@{}>{\raggedright\arraybackslash}p{4.3cm}X@{}}\toprule
Challenge & Solution \\\midrule
Broden server offline; candidate mirrors on Hugging Face were other data & Internet Archive copy, downloaded in parallel byte ranges; the archive index reported the compressed size, which first gave a truncated file --- detected with the SHA-1 digest, missing part fetched, file verified byte-identical \\
RG-encoded label PNGs, up to 5 overlapping label maps per pixel, image-level scene/texture labels & decoder merging all maps into $\leq$5 layers without counting a label twice per pixel; scene/texture as full-image masks. Our per-label image counts equal Broden's \texttt{label.csv} for every single-category material, scene and texture label \\
IoU only over images having the concept's category; labels in several categories & primary category per label; category-presence matrix per image; unit test against a brute-force IoU \\
$\approx10^{15}$ operations for a naive IoU & per-batch GPU \texttt{index\_add} into a concept$\times$unit table (exact counts) --- 2 min per layer \\
Exact quantile over up to $2\cdot10^8$ values per unit; our first version used $>$200\,GB RAM on Places365, where hundreds of units are zero on $>$99.5\% of positions (tie at 0) & redesigned: count values above/equal to candidate levels, detect ties explicitly ($T_k=0$), keep only values above a bounded level ($\leq$6\,GB); strict $S>T$ masks; unit test for ties \\
In-place ReLU overwrote tensors captured by hooks (breaks pre-ReLU / negative dissection) & ReLUs set to non-in-place, captured tensors cloned; unit test shows the problem \\
Anchoring upsampling at receptive-field centres & receptive-field start/stride computed from the layer geometry, sampled with \texttt{grid\_sample} \\
Fractional rotation $Q^\alpha$ & real Schur form, angles scaled by $\alpha$; orthogonality error $<10^{-14}$ \\
Fidelity test saturated (10-way accuracy 98\%) & 1000-way accuracy, grouped permutation, 10{,}000 random groups for the one pre-specified test \\\bottomrule
\end{tabularx}

\section{Limitations and conclusion}
Our AlexNet differs from the paper's Caffe model, the training experiments are at small scale (one run each for
dropout/BN/width), the ImageNet rotation curve has one seed, bootstrap intervals are approximate, and the human study
\fillin{was / was not} run. Network Dissection gives an objective, reproducible way to compare the interpretability of
representations, and its central finding --- interpretability is axis-aligned and not implied by discriminative power ---
holds in our implementation. But a label is a correlation, not an explanation of the model's decision: single labelled
units have little causal effect, labels depend on two hand-set thresholds and on how the dataset is labelled, and concepts
are better captured by directions than by units. It should therefore be used together with fidelity tests and other methods,
as Molnar recommends for any interpretation method.

\noindent\textbf{Declaration.} \fillin{State the use of AI tools as required by the course policy.}

\begin{thebibliography}{9}\small
\bibitem{bau2017} D. Bau, B. Zhou, A. Khosla, A. Oliva, A. Torralba. Network Dissection: Quantifying Interpretability of Deep Visual Representations. CVPR 2017. doi:10.1109/CVPR.2017.354
\bibitem{molnar} C. Molnar. \emph{Interpretable Machine Learning}. \url{https://christophm.github.io/interpretable-ml-book}, accessed \fillin{date}.
\bibitem{net2vec} R. Fong, A. Vedaldi. Net2Vec: Quantifying and Explaining how Concepts are Encoded by Filters in Deep Neural Networks. CVPR 2018.
\bibitem{adebayo} J. Adebayo et al. Sanity Checks for Saliency Maps. NeurIPS 2018.
\bibitem{fisher} A. Fisher, C. Rudin, F. Dominici. All Models are Wrong, but Many are Useful. JMLR 2019.
\bibitem{places} B. Zhou et al. Places: A 10 Million Image Database for Scene Recognition. TPAMI 2018.
\bibitem{imagenette} J. Howard. Imagenette. \url{https://github.com/fastai/imagenette}.
\end{thebibliography}
\end{document}
